%!TEX program = xelatex
\documentclass[lang=cn,11pt]{elegantbook}

% =============================================================================
% 核心宏包加载与版面紧凑化重载 (极致省纸打印、高清晰度、绝对连贯)
% =============================================================================
\usepackage{tabularx}   % 支持自适应折行表格
\usepackage{booktabs}   % 高级三线表，物理印刷质感极佳
\usepackage{enumitem}   % 精细化列表行距压缩
\usepackage{float}      % 控制表格/图片位置
\usepackage{amssymb}    % 提供数学符号 \blacktriangleright 支持

% 覆盖 ElegantBook 默认大页边距：将左右边距缩小为极高纸张利用率的 1.8cm，极大缩减打印总页数
\geometry{a4paper, left=1.8cm, right=1.8cm, top=2.2cm, bottom=2.2cm, footskip=0.8cm}

% 全局列表紧凑化设置，彻底节约纸张
\setlist{nosep, leftmargin=2em}

% =============================================================================
% 品牌与设计规范: 严格复古色系 (深墨绿 + 铜棕)
% =============================================================================
\definecolor{ecolor}{RGB}{46,76,56}        % 主色调: 复古深墨绿 (Pine Green) 用于标题及框架
\definecolor{accentcolor}{RGB}{143,91,46}  % 辅助色: 复古铜棕色 (Copper Brown) 用于词汇高亮与强调

% 核心高亮与样式宏定义
\newcommand{\retrokey}[1]{{\color{accentcolor}\textbf{#1}}}
\newcommand{\tips}[1]{{\color{ecolor}\textbf{#1}}}

% 解决 XeLaTeX 下音标字符在默认 Roman 字体中缺失的警告，切换为 Arial 渲染音标
\renewcommand{\citshape}{\kaishu\fontspec{Arial}}

% =============================================================================
% 文档元数据
% =============================================================================
\title{考研英语长难句 66 句精讲全书\\{\large ——— 复习背诵与打印专版 ———}}
\author{Harry}
\date{}

% =============================================================================
% 正文开始
% =============================================================================
\begin{document}

% 强制应用紧致行距 (1.15) 与极小段距 (1.5pt)，彻底消除松散空白，使排版行云流水
\linespread{1.15}\selectfont
\setlength{\parskip}{1.5pt}

\maketitle

\frontmatter
\tableofcontents

\mainmatter

% =============================================================================
\chapter{考研长难句句法奠基}
% =============================================================================

\begin{introduction}
  \item 谓语的分类大纲 (及物与不及物)
  \item 系动词的四种条件判定
  \item 基本语法名词认知与核心结构
\end{introduction}

在正式开启考研英语长难句的实战精读之前，我们必须对英语的**谓语分类**与**基本句型结构**有极其扎实的底层认知。谓语是任何句子的核心骨架。

\section{核心句法大纲：谓语的分类}

根据谓语动词是否具有实际语义以及其是否需要宾语，我们可以将谓语划分为三大支柱体系。

\begin{proposition}[谓语分类大纲与系动词条件]
  \textbf{一、有实义动词（Action Verbs）}
  \begin{itemize}
    \item \textbf{不及物动词 (Vi)} $\rightarrow$ \textbf{主 + 谓} 结构。动词本身即可表达完整含义，无需承受者（如 \textit{She smiled.}）。
    \item \textbf{及物动词 (Vt)} $\rightarrow$ 需要承受者方可完整：
      \begin{itemize}
        \item \textbf{主 + 谓 + 宾} 结构（如 \textit{We study English.}）。
        \item \textbf{主 + 谓 + 双宾} (人 + 物) $\rightarrow$ 常接 \textbf{to 人}（表方向）或 \textbf{for 人}（表受益）。
        \item \textbf{主 + 谓 + 宾补} (宾语是人+物时为双宾；若补充说明宾语则为宾补，如 \textit{make him happy}）。
      \end{itemize}
  \end{itemize}
  
  \textbf{二、无实义动词 $\rightarrow$ 系动词 (Link Verbs) $\rightarrow$ 主 + 系 + 表 结构}
  
  连系动词不表达具体动作，仅在逻辑上连接主语和表语（表语多为形容词或名词，用以表达状态或属性）。共有以下四类：
  \begin{enumerate}
    \item \textbf{be 动词} $\rightarrow$ 单独使用时译为“是”，是系统中最基础的系动词。
    \item \textbf{状态变化类} $\rightarrow$ \textbf{get}, \textbf{become}, \textbf{turn}, \textbf{go}, \textbf{grow} （译为：变得、成为）。
    \item \textbf{感官类} $\rightarrow$ \textbf{look}, \textbf{sound}, \textbf{smell}, \textbf{taste}, \textbf{feel} （译为：看/听/闻/尝/感觉起来）。
    \item \textbf{保持/证明类} $\rightarrow$ \textbf{seem}, \textbf{appear} (似乎、好像)；\textbf{keep}, \textbf{remain} (保持)；\textbf{prove} (证明、显示出)。
  \end{enumerate}
\end{proposition}

% =============================================================================
\chapter{第一辑：真题长难句精讲精练}
% =============================================================================

\begin{introduction}
  \item 第 1 句 $\rightarrow$ 英文报纸艺术板块报道的衰退 (1999-1)
  \item 第 2 句 $\rightarrow$ 大都市报纸文艺评论的消失 (1999-1)
  \item 第 3 句 $\rightarrow$ 早期社会仪式、神话与口头传统 (1999-5)
  \item 第 4 句 $\rightarrow$ 20世纪最重要的报纸评论文集 (1999-5)
  \item 第 5 句 $\rightarrow$ 二战前夕新闻纸张的廉价装饰 (1999-5)
  \item 第 6 句 $\rightarrow$ 纽曼对新闻写作的讽刺性定义 (1999-5)
  \item 第 7 句 $\rightarrow$ 商业方法专利对道富银行案的限制 (2000-1)
  \item 第 8 句 $\rightarrow$ 联邦巡回上诉法院重新审理裁决 (2000-1)
\end{introduction}

本辑精选了 1999 年至 2000 年考研英语真题中的 8 句经典长难句，进行深度的句法层级拆解与词汇自测。

\section{第 1 句：新闻艺术报道的衰退 (1999-1)}

\begin{note}
  Of all the changes that have taken place in English-language newspapers during the past quarter-century, perhaps the most far-reaching has been the inexorable decline in the scope and seriousness of their arts coverage. 【33 词】
\end{note}

\subsubsection*{$\blacktriangleright$ 句法层级拆解}
\noindent
\begin{tabularx}{\textwidth}{p{3.2cm} X X}
  \toprule
  \textbf{语法成分} & \textbf{具体英文内容} & \textbf{解析与修饰关系} \\
  \midrule
  \textbf{状语 (介词短语)} & Of all the changes & 在所有变化中（全句的限定范围状语） \\
  \midrule
  \textbf{定语从句} & that have taken place in English-language newspapers during the past quarter-century & 修饰 changes（that 引导定语从句，在从句中作主语，修饰变化发生的时间与范围） \\
  \midrule
  \textbf{主语 (倒装省略)} & the most far-reaching & 影响最深远的（省略了名词 change，起强调用法） \\
  \midrule
  \textbf{谓语 (系动词)} & has been & 已经是（现在完成时态） \\
  \midrule
  \textbf{表语/宾语} & the inexorable decline in the scope and seriousness of their arts coverage & 其艺术报道在覆盖范围和严肃程度上的不可阻挡的下滑 \\
  \bottomrule
\end{tabularx}

\begin{proposition}[倒装与语法细节辨析]
  \textbf{1. 倒装逻辑}：该句的正常语序为 \textit{the inexorable decline has been the most far-reaching (change)}。作者为了强调 \textbf{far-reaching}（影响深远）而将其前置，形成了主表倒装的艺术效果。\\
  \textbf{2. 常见不可数名词辨析 (information, data, arts, wood)}：
  \begin{itemize}
    \item \textbf{information}：绝对不可数，表示抽象的“信息”。
    \item \textbf{data}：在现代正式学术文体中作不可数名词，表示“数据、资料”。
    \item \textbf{arts}：当特指具体的“艺术门类、人文学科”时常以复数形式出现（如 \textit{liberal arts} 文科）。
    \item \textbf{wood}：不可数表示“木材”；当其具有明确物理边界时，使用复数形式 \textbf{woods}（表示树林）。
  \end{itemize}
\end{proposition}

\subsubsection*{$\blacktriangleright$ 重点词汇与音标}
\noindent
\begin{tabularx}{\textwidth}{l X p{7cm}}
  \toprule
  \textbf{单词} & \textbf{音标与词性} & \textbf{核心释义} \\
  \midrule
  \retrokey{change} & {\citshape /tʃeɪndʒ/} \textit{n./v.} & 变化；改变，转换 \\
  \midrule
  \retrokey{take place} & {\citshape /teɪk pleɪs/} \textit{phr.} & 发生，举行（注意：\textbf{绝对没有被动语态}） \\
  \midrule
  \retrokey{far-reaching} & {\citshape /fɑː ˈriːtʃɪŋ/} \textit{adj.} & 影响深远的，波及广泛的 \\
  \midrule
  \retrokey{inexorable} & {\citshape /ɪnˈeksərəbl/} \textit{adj.} & 不可阻挡的，无情的，无法改变的 \\
  \midrule
  \retrokey{decline} & {\citshape /dɪˈklaɪn/} \textit{n./v.} & 下降，衰退，减少，谢绝 \\
  \midrule
  \retrokey{scope} & {\citshape /skəʊp/} \textit{n.} & 范围，领域，眼界 \\
  \midrule
  \retrokey{seriousness} & {\citshape /ˈsɪəriəsnəs/} \textit{n.} & 严肃性，严肃程度，严重性 \\
  \midrule
  \retrokey{coverage} & {\citshape /ˈkʌvərɪdʒ/} \textit{n.} & 新闻报道，覆盖范围（媒体） \\
  \bottomrule
\end{tabularx}

\subsubsection*{$\blacktriangleright$ 翻译与真题启示}
\begin{remark}
  \textbf{参考译文}：在过去25年里，英文报纸上所发生的所有变化中，影响最为深远的也许是艺术板块报道的覆盖范围和严肃程度的不断下滑。
\end{remark}

\par\vspace{10pt}\hrule\vspace{10pt}

% =============================================================================
\section{第 2 句：大都市报纸中的文艺评论 (1999-1)}

\begin{note}
  It is difficult to the point of impossibility for the average reader under the age of forty to imagine a time when high-quality arts criticism could be found in most big-city newspapers. 【32 词】
\end{note}

\subsubsection*{$\blacktriangleright$ 句法层级拆解}
\noindent
\begin{tabularx}{\textwidth}{p{3.2cm} X X}
  \toprule
  \textbf{语法成分} & \textbf{具体英文内容} & \textbf{解析与修饰关系} \\
  \midrule
  \textbf{形式主语} & It & 代指后文的不定式，无实际语义 \\
  \midrule
  \textbf{谓语与表语} & is difficult to the point of impossibility & 困难到了几乎不可能的程度（difficult 为表语，后接程度修饰短语） \\
  \midrule
  \textbf{逻辑主语} & for the average reader under the age of forty & 对四十岁以下的普通读者而言（\textbf{for sb.} 结构作不定式的逻辑主语） \\
  \midrule
  \textbf{真正主语 (不定式)} & to imagine a time when… & 想象曾经有一个时代（真正的主语） \\
  \midrule
  \textbf{定语从句} & when high-quality arts criticism could be found in most big-city newspapers & 修饰时间名词 time（when 引导定语从句，在从句中作时间状语） \\
  \bottomrule
\end{tabularx}

\begin{proposition}[前缀衍生规律与被动译主动]
  \textbf{1. 否定前缀 im- / in- 规律}：
  \begin{itemize}
    \item \textbf{表示否定}：\textit{impossible}（不可能的）、\textit{incomplete}（不完全的）、\textit{illegal}（非法的，il-前缀）。
    \item \textbf{表示方向（向内）}：\textit{import}（进口）、\textit{inspect}（视察，往里看）、\textit{impose}（强加，im + pose 放 $\rightarrow$ 强加于上）。
  \end{itemize}
  \textbf{2. 动词化前缀 em- / en- (使动化)}：\\
  在名词或形容词前加上 \textbf{em- / en-}，可使其转换为及物动词：
  \begin{itemize}
    \item \textbf{enlarge} (\textit{large} 大 $\rightarrow$ 使扩大) \quad \textbf{embody} (\textit{body} 身体 $\rightarrow$ 使体现/具体化)
    \item \textbf{widen} (\textit{wide} 宽 $\rightarrow$ 使加宽) \quad \textbf{lengthen} (\textit{length} 长度 $\rightarrow$ 使延长)
    \item \textbf{strengthen} (\textit{strength} 力量 $\rightarrow$ 使加强) \quad \textbf{sharpen} (\textit{sharp} 锋利 $\rightarrow$ 使锋利/削尖)
  \end{itemize}
  \textbf{3. 翻译技巧：被动句 $\rightarrow$ 主动译}：原文中的被动语态 \textbf{could be found} 在中文表达中应化被动为主动，译为“……人们可以从大多数城市报纸上找到……”，更加顺畅自然。
\end{proposition}

\subsubsection*{$\blacktriangleright$ 重点词汇与音标}
\noindent
\begin{tabularx}{\textwidth}{l X p{7cm}}
  \toprule
  \textbf{单词} & \textbf{音标与词性} & \textbf{核心释义} \\
  \midrule
  \retrokey{point} & {\citshape /pɔɪnt/} \textit{n.} & 点，观点，时刻；程度（\textbf{to the point of} = 到了……的程度） \\
  \midrule
  \retrokey{average} & {\citshape /ˈævərɪdʒ/} \textit{adj./n.} & 普通的，平均的；平均数（区别于 \textbf{ordinary} 平凡的） \\
  \midrule
  \retrokey{impossibility} & {\citshape /ɪmˌpɒsəˈbɪləti/} \textit{n.} & 不可能，不可能的事（\textit{im-} 否定前缀） \\
  \midrule
  \retrokey{imagine} & {\citshape /ɪˈmædʒɪn/} \textit{v.} & 想象，设想，认为 \\
  \midrule
  \retrokey{envision} & {\citshape /ɪnˈvɪʒn/} \textit{v.} & 想象，展望，预想（比 imagine 更加正式） \\
  \midrule
  \retrokey{criticism} & {\citshape /ˈkrɪtɪsɪzəm/} \textit{n.} & 批评，批判，文艺评论，评论文章 \\
  \bottomrule
\end{tabularx}

\subsubsection*{$\blacktriangleright$ 翻译与真题启示}
\begin{remark}
  \textbf{参考译文}：对于大多数四十岁以下的普通读者来说，很难想象（甚至几乎不可能想象）在大多数大城市报纸上曾经能找到高质量的艺术评论的那个时代。
\end{remark}

\par\vspace{10pt}\hrule\vspace{10pt}

% =============================================================================
\section{第 3 句：早期社会宗教仪式与神话 (1999-5)}

\begin{note}
  For example, some early societies ceased to consider certain rites essential to their well-being and abandoned them; nevertheless, they retained as parts of their oral tradition the myths that had grown up around the rites and admired them for their artistic qualities rather than for their religious usefulness. 【48 词】
\end{note}

\subsubsection*{$\blacktriangleright$ 句法层级拆解}
\noindent
\begin{tabularx}{\textwidth}{p{3.2cm} X X}
  \toprule
  \textbf{语法成分} & \textbf{具体英文内容} & \textbf{解析与修饰关系} \\
  \midrule
  \textbf{第一主句} & some early societies ceased to consider certain rites essential to their well-being and abandoned them & 早期社会不再视特定仪式为其福祉所必需，从而抛弃了它们（\textbf{consider A B} 结构，B为形容词） \\
  \midrule
  \textbf{转折连词} & nevertheless & 然而，尽管如此（连接两个语义并列且高度对立的句子） \\
  \midrule
  \textbf{第二主句主干} & they retained … the myths … and admired them … & 他们保留了神话并欣赏它们（宾语 details 在后文） \\
  \midrule
  \textbf{宾语前置状语} & as parts of their oral tradition & 作为口头传统的一部分（原本在 myths 后面，被强制前置，增加了句法难度） \\
  \midrule
  \textbf{定语从句} & that had grown up around the rites & 修饰 myths（that 引导定语从句在从句中作主语） \\
  \midrule
  \textbf{介词状语} & for their artistic qualities rather than for their religious usefulness & 因为它们的艺术价值而非宗教功用（由 rather than 连接的两个并列原因状语） \\
  \bottomrule
\end{tabularx}

\begin{proposition}[“把 A 看作 B”的经典句型拓展]
  在考研英语中，“把 A 看作 B”是高频出现的重点逻辑结构，包含以下主动与被动形式：
  \begin{table}[H]
    \centering
    \small
    \begin{tabularx}{\textwidth}{p{5.5cm} X}
      \toprule
      \textbf{主动表达形式} & \textbf{对应被动表达形式} \\
      \midrule
      \textbf{see A as B} & \textbf{A be seen as B} \\
      \textbf{view A as B} & \textbf{A be viewed as B} \\
      \textbf{consider A (as) B} & \textbf{A be considered (as/to be) B} \\
      \textbf{regard A as B} & \textbf{A be regarded as B} \\
      \textbf{think of A as B} & \textbf{A be thought of as B} \\
      \textbf{refer to A as B} & \textbf{A be referred to as B} \\
      \bottomrule
    \end{tabularx}
  \end{table}
  \textbf{语法补充}：\textbf{cease to do} 表示“停止做某事”（等同于 \textit{stop doing}）。\textbf{nevertheless} 作连词表强烈转折，语气比 \textit{however} 更加正式，通常放置于句首。
\end{proposition}

\subsubsection*{$\blacktriangleright$ 重点词汇与音标}
\noindent
\begin{tabularx}{\textwidth}{l X p{7cm}}
  \toprule
  \textbf{单词} & \textbf{音标与词性} & \textbf{核心释义} \\
  \midrule
  \retrokey{consider} & {\citshape /kənˈsɪdə/} \textit{v.} & 认为，考虑，体贴 \\
  \midrule
  \retrokey{considerable} & {\citshape /kənˈsɪdərəbl/} \textit{adj.} & 相当大的，大量的，值得考虑的 \\
  \midrule
  \retrokey{well-being} & {\citshape /wel ˈbiːɪŋ/} \textit{n.} & 健康，幸福，福祉（健康与幸福状态） \\
  \midrule
  \retrokey{myth} & {\citshape /mɪθ/} \textit{n.} & 神话故事，荒诞说法，错误观念 \\
  \midrule
  \retrokey{nevertheless} & {\citshape /ˌnevəðəˈles/} \textit{adv./conj.} & 然而，尽管如此，但是 \\
  \midrule
  \retrokey{rite} & {\citshape /raɪt/} \textit{n.} & 宗教仪式，传统典礼 \\
  \bottomrule
\end{tabularx}

\subsubsection*{$\blacktriangleright$ 翻译与真题启示}
\begin{remark}
  \textbf{参考译文}：例如，一些早期社会不再认为某些特定仪式对其幸福至关重要，因而抛弃了这些仪式；尽管如此，他们仍将围绕这些仪式发展起来的神话故事作为口头传统的一部分保留下来，并因其艺术价值而非宗教功用而推崇它们。
\end{remark}

\par\vspace{10pt}\hrule\vspace{10pt}

% =============================================================================
\section{第 4 句：20世纪最重要的报纸评论文集 (1999-5)}

\begin{note}
  Yet a considerable number of the most significant collections of criticism published in the 20th century consisted in large part of newspaper reviews. 【23 词】
\end{note}

\subsubsection*{$\blacktriangleright$ 句法层级拆解}
\noindent
\begin{tabularx}{\textwidth}{p{3.2cm} X X}
  \toprule
  \textbf{语法成分} & \textbf{具体英文内容} & \textbf{解析与修饰关系} \\
  \midrule
  \textbf{转折连词} & Yet & 然而（置于句首强转折） \\
  \midrule
  \textbf{主语 (核心)} & a considerable number of the most significant collections of criticism & 20世纪最重要的一部分评论文集（collections 为主语核心） \\
  \midrule
  \textbf{后置定语 (done)} & published in the 20th century & 在20世纪出版的（过去分词短语作后置定语，修饰 collections，等价于定语从句 \textit{which were published...}） \\
  \midrule
  \textbf{谓语} & consisted in large part of & 在很大程度上是由……构成的（consist of 的完成形式） \\
  \midrule
  \textbf{宾语} & newspaper reviews & 报纸评论 \\
  \bottomrule
\end{tabularx}

\subsubsection*{$\blacktriangleright$ 重点词汇与音标}
\begin{itemize}
  \item \retrokey{published} {\citshape /ˈpʌblɪʃt/} \textit{adj./done.} 已出版的，发表的（过去分词作定语表被动或完成）。
  \item \retrokey{in large part} {\citshape /ɪn lɑːdʒ pɑːt/} \textit{phr.} 在很大程度上，多半（等价于 \textit{largely, to a large extent}）。
  \item \retrokey{yet} {\citshape /jet/} \textit{adv./conj.} 然而，但是（表转折）；在否定句中译为“还，尚未”。
  \item \retrokey{review} {\citshape /rɪˈvjuː/} \textit{n./v.} 评论（critique），评价，回顾，复习。
\end{itemize}

\subsubsection*{$\blacktriangleright$ 翻译与真题启示}
\begin{remark}
  \textbf{参考译文}：然而，20世纪出版的大量最重要的评论文集在很大程度上由报纸评论构成。（或精简译法：然而，20世纪出版的大量最重要的评论文集多是报纸评论。）
\end{remark}

\par\vspace{10pt}\hrule\vspace{10pt}

% =============================================================================
\section{第 5 句：二战前夕新闻纸张的廉价装饰 (1999-5)}

\begin{note}
  We are even farther removed from the unfocused newspaper reviews published in England between the turn of the 20th century and the eve of World War II, at a time when newsprint was dirt-cheap and stylish arts criticism was considered an ornament to the publications in which it appeared. 【49 词】
\end{note}

\subsubsection*{$\blacktriangleright$ 句法层级拆解}
\noindent
\begin{tabularx}{\textwidth}{p{3.2cm} X X}
  \toprule
  \textbf{语法成分} & \textbf{具体英文内容} & \textbf{解析与修饰关系} \\
  \midrule
  \textbf{主句主干} & We are even farther removed from the unfocused newspaper reviews & 我们与那些散漫的报纸评论相去甚远（主系表 + 远离结构） \\
  \midrule
  \textbf{分词定语} & published in England between … and … & 在英国出版的，时间在……与……之间（过去分词短语作后置定语） \\
  \midrule
  \textbf{时间状语} & at a time when… & 在曾经的一个时代（when 引导定语从句修饰 a time） \\
  \midrule
  \textbf{when 从句主干} & newsprint was dirt-cheap and stylish arts criticism was considered an ornament… & 新闻纸张价格极其低廉，且新颖的艺术评论被视为一种装饰 \\
  \midrule
  \textbf{定语从句} & in which it appeared & 修饰 publications（介词 + 关系代词 引导定语从句，it 指代 criticism） \\
  \bottomrule
\end{tabularx}

\begin{proposition}[定语从句关系副词与汉语翻译增词法]
  \textbf{1. 关系副词 (where / why / when) 引导的定语从句判定}：
  \begin{itemize}
    \item \textbf{where} $\rightarrow$ 先行词为表示地点的名词（如 \textit{the place where I was born}）。
    \item \textbf{when} $\rightarrow$ 先行词为表示时间的名词（如本句中的 \textit{a time when...}）。
    \item \textbf{why} $\rightarrow$ 先行词通常只能是 \textit{reason}（如 \textit{the reason why he left}）。
  \end{itemize}
  \textbf{2. 翻译小技巧：汉语增词法}：\\
  汉语习惯于多用动词、副词和关联词。在将英语长句汉译时，应根据汉语表达习惯适当增加解释性字词，以使译文流畅：
  \begin{itemize}
    \item \textbf{增加过程词}：\textit{in evolution} $\rightarrow$ 在进化\retrokey{的过程}中。
    \item \textbf{增加数量词}：\textit{Flowers bloom} $\rightarrow$ \retrokey{无数}花朵盛开。
    \item \textbf{增加转折副词}：原文两句并列，译文可增“\retrokey{第一}……，\retrokey{第二}……”进行结构分层。
  \end{itemize}
\end{proposition}

\subsubsection*{$\blacktriangleright$ 重点词汇与音标}
\noindent
\begin{tabularx}{\textwidth}{l X p{7cm}}
  \toprule
  \textbf{单词} & \textbf{音标与词性} & \textbf{核心释义} \\
  \midrule
  \retrokey{newsprint} & {\citshape /ˈnjuːzprɪnt/} \textit{n.} & 新闻纸，低克重印刷纸 \\
  \midrule
  \retrokey{dirt-cheap} & {\citshape /ˌdɜːt ˈtʃiːp/} \textit{adj./adv.} & 极其便宜的，贱价的 \\
  \midrule
  \retrokey{ornament} & {\citshape /ˈɔːnəmənt/} \textit{n./v.} & 装饰品，点缀；装饰，美化 \\
  \midrule
  \retrokey{unfocused} & {\citshape /ʌnˈfəʊkəst/} \textit{adj.} & 漫无边际的，焦点模糊的 \\
  \bottomrule
\end{tabularx}

\subsubsection*{$\blacktriangleright$ 翻译与真题启示}
\begin{remark}
  \textbf{参考译文}：我们与那个时代相去甚远——那是20世纪初至二战前夕，在英国出版的漫无边际的报纸评论的时代，那时新闻纸价格极其低廉，新颖的艺术评论被视为其所在出版物的一种装饰。
\end{remark}

\par\vspace{10pt}\hrule\vspace{10pt}

% =============================================================================
\section{第 6 句：纽曼对新闻写作的讽刺性定义 (1999-5)}

\begin{note}
  "So few authors have brains enough or literary gift enough to keep their own end up in journalism," Newman wrote; "that I am tempted to define 'journalism' as 'a term of contempt applied by writers who are not read to writers who are.'" 【43 词】
\end{note}

\subsubsection*{$\blacktriangleright$ 句法层级拆解}
\noindent
\begin{tabularx}{\textwidth}{p{3.2cm} X X}
  \toprule
  \textbf{语法成分} & \textbf{具体英文内容} & \textbf{解析与修饰关系} \\
  \midrule
  \textbf{主句主干 (前)} & So few authors have brains enough ... to keep their own end up in journalism & 拥有足够智慧在新闻界维护自身立场的作家如此之少（so... that 倒装并列结构） \\
  \midrule
  \textbf{插入语} & Newman wrote & 纽曼写道（置于分号中间作插入） \\
  \midrule
  \textbf{主句主干 (后)} & that I am tempted to define 'journalism' as 'a term of contempt...' & 以至于我忍不住把“新闻写作”定义为一种蔑视之辞（that 引导结果状语从句） \\
  \midrule
  \textbf{定语修饰} & applied by writers who are not read to writers who are (read) & 被不受欢迎的作家施加于受欢迎的作家之上的（过去分词短语作后置定语，修饰 term） \\
  \midrule
  \textbf{定语从句与省略} & who are not read ... who are (read) & 修饰 writers（两个定语从句，最后一个省略了 read，表被动） \\
  \bottomrule
\end{tabularx}

\begin{proposition}[定语从句的简化与 There is 结构的特殊情况]
  \textbf{1. 定语从句的简化法则 (从句 $\rightarrow$ 短语)}：\\
  当定语从句的关系代词充当从句的**主语**时，可直接将其简化为分词短语（主动关系用 \textbf{doing}，被动关系用 \textbf{done}）：
  \begin{itemize}
    \item \textit{which belongs to...} $\rightarrow$ \textbf{belonging to...}
    \item \textit{which are written in...} $\rightarrow$ \textbf{written in...}
    \item \textit{who is standing at...} $\rightarrow$ \textbf{standing at...}
  \end{itemize}
  \textbf{2. There is 结构的三大特殊阅读陷阱}：
  \begin{enumerate}
    \item \textbf{There is + 主语 + 目的状语} $\rightarrow$ 表示在某处有某事物去执行某个目的。
    \item \textbf{介词强制前置} $\rightarrow$ 造成主谓宾的逻辑分割。
    \item \textbf{后接不定式作目的/结果状语} $\rightarrow$ 如 \textit{ask sb. to do} 等，阅读时需快速提炼动宾主干。
  \end{enumerate}
\end{proposition}

\subsubsection*{$\blacktriangleright$ 重点词汇与音标}
\noindent
\begin{tabularx}{\textwidth}{l X p{7cm}}
  \toprule
  \textbf{单词} & \textbf{音标与词性} & \textbf{核心释义} \\
  \midrule
  \retrokey{author} & {\citshape /ˈɔːθə/} \textit{n.} & 作者，发起人（词根 \textit{auth-} 表示增加） \\
  \midrule
  \retrokey{term} & {\citshape /tɜːm/} \textit{n./v.} & 术语，学期，条款；把……称为 \\
  \midrule
  \retrokey{keep one's end up} & {\citshape /kiːp wʌnz end ʌp/} \textit{phr.} & 坚持到底，在困难中维护自己的立场 \\
  \midrule
  \retrokey{augment} & {\citshape /ɔːɡˈment/} \textit{v.} & 增大，扩张，增加（如增加收入） \\
  \midrule
  \retrokey{auction} & {\citshape /ˈɔːkʃn/} \textit{n./v.} & 拍卖 \\
  \bottomrule
\end{tabularx}

\subsubsection*{$\blacktriangleright$ 翻译与真题启示}
\begin{remark}
  \textbf{参考译文}：纽曼写道：“有足够智慧或文学才华坚持在新闻界闯荡的作家少之又少，以至于我忍不住把‘新闻写作’定义为‘不受读者欢迎的作家’施用于‘受读者欢迎的作家’身上的蔑视之辞。”
\end{remark}

\par\vspace{10pt}\hrule\vspace{10pt}

% =============================================================================
\section{第 7 句：商业方法专利对道富银行案的限制 (2000-1)}

\begin{note}
  Curbs on business-method claims would be a dramatic about-face, because it was the Federal Circuit itself that introduced such patents with its 1998 decision in the so-called State Street Bank case, approving a patent on a way of pooling mutual-fund assets. 【41 词】
\end{note}

\subsubsection*{$\blacktriangleright$ 句法层级拆解}
\noindent
\begin{tabularx}{\textwidth}{p{3.2cm} X X}
  \toprule
  \textbf{语法成分} & \textbf{具体英文内容} & \textbf{解析与修饰关系} \\
  \midrule
  \textbf{主句主干} & Curbs on business-method claims would be a dramatic about-face & 对商业方法专利申请的限制将是一次重大的逆转（主系表结构） \\
  \midrule
  \textbf{原因状语从句} & because it was the Federal Circuit itself that introduced such patents… & 因为正是联邦巡回法院本身引入了这类专利（because 引导原因状语从句） \\
  \midrule
  \textbf{强调句结构} & \textbf{it was … that} & 强调原因从句的主语 the Federal Circuit itself \\
  \midrule
  \textbf{伴随状语 (分词)} & approving a patent on a way of pooling mutual-fund assets & 同时批准了一项关于集中基金资产方式的专利（现在分词短语作伴随状语） \\
  \bottomrule
\end{tabularx}

\begin{proposition}[强调句结构的精准判定]
  \textbf{强调句的核心逻辑式}：\textbf{It is/was + 被强调部分 + that/who + 其余部分}。\\
  \textbf{工程化检验方法}：若把 \textit{it is/was ... that} 强制从句中删去，剩下的内容如果仍然是一个**结构完整、语义通顺的完整句子**，则该句必定是强调句。
  \begin{itemize}
    \item \textbf{强调时间} $\rightarrow$ \textit{It was in 1998 that...} (正是于1998年……)
    \item \textbf{强调地点} $\rightarrow$ \textit{It was here that...} (正是在这里……)
  \end{itemize}
  \textbf{语法细节}：末尾的 \textbf{approving} 分词短语充当伴随状语，表示该动作与主句谓语同时发生，翻译时可译为“同时/从而……”。
\end{proposition}

\subsubsection*{$\blacktriangleright$ 重点词汇与音标}
\noindent
\begin{tabularx}{\textwidth}{l X p{7cm}}
  \toprule
  \textbf{单词} & \textbf{音标与词性} & \textbf{核心释义} \\
  \midrule
  \retrokey{curb} & {\citshape /kɜːb/} \textit{n./v.} & 限制，管制，控制（\textbf{curb on} 对……的限制） \\
  \midrule
  \retrokey{about-face} & {\citshape /əˌbaʊt ˈfeɪs/} \textit{n.} & 180度大转变，戏剧性的大逆转 \\
  \midrule
  \retrokey{patent} & {\citshape /ˈpeɪtnt/} \textit{n./v.} & 专利，专利权；申请专利 \\
  \midrule
  \retrokey{pool} & {\citshape /puːl/} \textit{v./n.} & 集中，汇聚（如汇聚资产）；水池，联营 \\
  \midrule
  \retrokey{mutual-fund} & {\citshape /ˌmjuːtʃuəl ˈfʌnd/} \textit{n.} & 共同基金 \\
  \bottomrule
\end{tabularx}

\subsubsection*{$\blacktriangleright$ 翻译与真题启示}
\begin{remark}
  \textbf{参考译文}：对商业方法专利申请的限制将是一次重大转变，因为正是联邦巡回上诉法院本身通过其1998年在所谓“美国道富银行案”中的裁决引入了这类专利，同时批准了一项关于集中管理共同基金资产方式的专利。
\end{remark}

\par\vspace{10pt}\hrule\vspace{10pt}

% =============================================================================
\section{第 8 句：联邦巡回上诉法院的命令 (2000-1)}

\begin{note}
  The Federal Circuit issued an unusual order, stating that the case would be heard by all 12 of the court's judges, rather than a typical panel of three, and that one issue it wants to evaluate is whether it should "reconsider" its State Street Bank ruling. 【46 词】
\end{note}

\subsubsection*{$\blacktriangleright$ 句法层级拆解}
\noindent
\begin{tabularx}{\textwidth}{p{3.2cm} X X}
  \toprule
  \textbf{语法成分} & \textbf{具体英文内容} & \textbf{解析与修饰关系} \\
  \midrule
  \textbf{主句主干} & The Federal Circuit issued an unusual order & 联邦巡回法院发布了一项不同寻常的命令（主谓宾结构） \\
  \midrule
  \textbf{伴随状语 (分词)} & stating that… & 声明说……（现在分词短语作伴随状语，具体阐述 order 的内容，等价于定语从句 \textit{which states that...}） \\
  \midrule
  \textbf{并列宾语从句} & (1) that the case would be heard by all 12 judges… and (2) that one issue … is whether… & 声明内容之一：全院12名法官审理；声明内容之二：评估是否重新考虑裁决（两个 that 从句并列作 stating 的宾语） \\
  \midrule
  \textbf{表语从句} & whether it should "reconsider" its State Street Bank ruling & 是否应该重新考虑道富银行案的裁决（whether 引导表语从句） \\
  \bottomrule
\end{tabularx}

\begin{proposition}[分词短语降维与代词链还原]
  \textbf{1. 分词短语降维}：在高级学术英语中，常使用 \textbf{stating that...} 代替复杂的定语从句 \textit{which states that...}。这是为了减少从句的多层嵌套，使句子结构更加紧凑利落。\\
  \textbf{2. 翻译的增词逻辑}：面对 A and B 的长并列宾语从句，直接翻译容易造成句子臃肿。在翻译时，可以采用“**第一**……，**第二**……”的分层增词法，使结构层次清晰。\\
  \textbf{3. 代词指代还原 (代词链)}：
  \begin{itemize}
    \item 第一个 \textit{it} (it wants to evaluate) $\rightarrow$ 还原指代主语 \textbf{the Federal Circuit}（联邦巡回法院）。
    \item 第二个 \textit{it} (it should "reconsider") $\rightarrow$ 同样还原指代 \textbf{the Federal Circuit}。
  \end{itemize}
\end{proposition}

\subsubsection*{$\blacktriangleright$ 重点词汇与音标}
\noindent
\begin{tabularx}{\textwidth}{l X p{7cm}}
  \toprule
  \textbf{单词} & \textbf{音标与词性} & \textbf{核心释义} \\
  \midrule
  \retrokey{disorder} & {\citshape /dɪsˈɔːdə/} \textit{n.} & 混乱，失调，骚乱（\textit{dis-} 否定前缀） \\
  \midrule
  \retrokey{panel} & {\citshape /ˈpænl/} \textit{n.} & 陪审团，专家小组，面板，讨论小组 \\
  \midrule
  \retrokey{issue} & {\citshape /ˈɪʃuː/} \textit{n./v.} & 问题，争端；期（期刊）；发布，发行 \\
  \midrule
  \retrokey{ruling} & {\citshape /ˈruːlɪŋ/} \textit{n./adj.} & （法庭的）裁决，裁定；统治的，支配的 \\
  \bottomrule
\end{tabularx}

\subsubsection*{$\blacktriangleright$ 翻译与真题启示}
\begin{remark}
  \textbf{参考译文}：联邦巡回法院发布了一项不同寻常的命令，该命令声明了两点：第一，该案不应由典型的三人法官小组审理，而应由法院全部12名法官共同审听；第二，联邦巡回法院希望衡量是否应该“重新考虑”美国道富银行案的裁决。
\end{remark}

\par\vspace{10pt}\hrule\vspace{15pt}

% =============================================================================
% 第三部分：真题长难句实战进阶
% =============================================================================
\chapter{第二辑：真题长难句实战进阶}
% =============================================================================

\begin{introduction}
  \item 第 9 句 $\rightarrow$ 引爆点中的社会流行元素与极少数人行为 (2010-1)
  \item 第 10 句 $\rightarrow$ 社会影响力与名人的媒体人际效应 (2010-1)
  \item 第 11 句 $\rightarrow$ 会计准则对第三方资产评估的规定 (2010-4)
  \item 第 12 句 $\rightarrow$ 国际会计标准委员会规则重组压力 (2010-4)
  \item 第 13 句 $\rightarrow$ 麦克里维警告IASB政治现实真空 (2010-4)
  \item 第 14 句 $\rightarrow$ 萧条市场与银行因恐亏损瘫痪 (2010-4)
  \item 第 15 句 $\rightarrow$ FASB与IASB不顾集团反对整顿期权 (2010-4)
  \item 第 16 句 $\rightarrow$ 丰富多样刺激因素环境对智力的开发 (1990-2)
\end{introduction}

本辑精选了 2010 年考研真题（及 1990 年经典长难句）中极具学术厚度与句法难度的 8 个长难句，包含大量嵌套定语从句和并列省略考点。

\section{第 9 句：引爆点中的社会流行元素 (2010-1)}

\begin{note}
  In his book The Tipping Point, Malcolm Gladwell argues that "social epidemics" are driven in large part by the action of a tiny minority of special individuals, often called influentials, who are unusually informed, persuasive, or well-connected. 【37 词】
\end{note}

\subsubsection*{$\blacktriangleright$ 句法层级拆解}
\noindent
\begin{tabularx}{\textwidth}{p{3.2cm} X X}
  \toprule
  \textbf{语法成分} & \textbf{具体英文内容} & \textbf{解析与修饰关系} \\
  \midrule
  \textbf{状语 (介词短语)} & In his book The Tipping Point & 在其著作《引爆点》中（全句的限定状语） \\
  \midrule
  \textbf{主语与谓语} & Malcolm Gladwell argues & 马尔科姆·格拉德威尔指出（主谓结构） \\
  \midrule
  \textbf{宾语从句} & that "social epidemics" are driven ... by the action of a tiny minority of special individuals & 宾语从句。social epidemics 为从句主语，整体为被动语态 \\
  \midrule
  \textbf{后置定语 (done)} & often called influentials & 通常被称为有影响力人士（分词短语作后置同位语，修饰 individuals，简化自定语从句） \\
  \midrule
  \textbf{定语从句} & who are unusually informed, persuasive, or well-connected & 修饰 individuals（who 引导定语从句，与分词短语构成双重定语修饰） \\
  \bottomrule
\end{tabularx}

\begin{proposition}[定语从句多重简化的底层逻辑]
  \textbf{关系代词作主语的定语从句简化法}：
  \begin{itemize}
    \item \textbf{A: be动词 + 形容词短语} $\rightarrow$ 直接删除关系词与 be 动词，将形容词短语后置修饰。
    \item \textbf{B: be动词 + 过去分词 (被动)} $\rightarrow$ 直接简化为 \textbf{done} 过去分词短语（如本句中的 \textit{who are often called} $\rightarrow$ \textbf{often called}）。
    \item \textbf{C: be动词 + 现在分词 (主动)} $\rightarrow$ 直接简化为 \textbf{doing} 分词短语。
  \end{itemize}
  \textbf{翻译技巧}：对于被动语态 \textbf{are driven by...}，在汉译时常采用“……由……所引领/驱动”的主动句式表达，更为符合中文逻辑。
\end{proposition}

\subsubsection*{$\blacktriangleright$ 重点词汇与音标}
\noindent
\begin{tabularx}{\textwidth}{l X p{7cm}}
  \toprule
  \textbf{单词} & \textbf{音标与词性} & \textbf{核心释义} \\
  \midrule
  \retrokey{tip} & {\citshape /tɪp/} \textit{n./v.} & 顶端，小费，提示；倾斜（\textbf{tipping point} = 临界点，引爆点） \\
  \midrule
  \retrokey{epidemic} & {\citshape /ˌepɪˈdemɪk/} \textit{n./adj.} & 流行病，蔓延；流行性的（\textit{demo-} 词根表示人民） \\
  \midrule
  \retrokey{democratic} & {\citshape /ˌdeməˈkrætɪk/} \textit{adj.} & 民主的，民主政治的 \\
  \midrule
  \retrokey{influential} & {\citshape /ˌɪnfluˈenʃl/} \textit{adj./n.} & 有影响力的；有影响力的人物（复数形式 \textbf{influentials}） \\
  \bottomrule
\end{tabularx}

\subsubsection*{$\blacktriangleright$ 翻译与真题启示}
\begin{remark}
  \textbf{参考译文}：马尔科姆·格拉德威尔在其《引爆点》一书中指出，“社会上的流行元素”在很大程度上是由极少数特殊群体的行为所引领的。这些人通常被称为“有影响力人士”，他们极其见多识广，善于言辞或善于交际。
\end{remark}

\par\vspace{10pt}\hrule\vspace{10pt}

% =============================================================================
\section{第 10 句：社会影响与名人的媒体效应 (2010-1)}

\begin{note}
  The researchers' argument stems from a simple observation about social influence: with the exception of a few celebrities like Oprah Winfrey—whose outsized presence is primarily a function of media, not interpersonal, influence—even the most influential members of a population simply don't interact with that many others. 【48 词】
\end{note}

\subsubsection*{$\blacktriangleright$ 句法层级拆解}
\noindent
\begin{tabularx}{\textwidth}{p{3.2cm} X X}
  \toprule
  \textbf{语法成分} & \textbf{具体英文内容} & \textbf{解析与修饰关系} \\
  \midrule
  \textbf{主句主干} & The researchers' argument stems from a simple observation about social influence & 研究者们的观点源于对社会影响力的一次简单调查（主谓宾结构） \\
  \midrule
  \textbf{冒号后展开} & with the exception of … even the most influential members … don't interact … & 冒号后具体分述：除了少数名人，即使最核心的群体也根本不与那么多人互动 \\
  \midrule
  \textbf{破折号定语从句} & whose outsized presence is primarily a function of media, not interpersonal, influence & 插入修饰 celebrities（whose 引导非限制性定语从句，作破折号插入语） \\
  \midrule
  \textbf{核心分述主干} & even the most influential members of a population simply don't interact with that many others & 即使是人群中最有影响力的个体也根本不会与那么多人发生互动 \\
  \bottomrule
\end{tabularx}

\begin{proposition}[冒号指示逻辑与 few 词义辨析]
  \textbf{1. 标点符号的解题预判 $\rightarrow$ 冒号}：冒号在阅读中表示“总述 $\rightarrow$ 分述”或“抽象 $\rightarrow$ 具体”。冒号前面的内容是抽象的，冒号后的分述主干才是命题的核心所在。\\
  \textbf{2. few 词族高频辨析}：
  \begin{table}[H]
    \centering
    \small
    \begin{tabularx}{\textwidth}{p{3cm} X X}
      \toprule
      \textbf{表达形式} & \textbf{修饰的名词类型} & \textbf{核心语义 (肯定/否定)} \\
      \midrule
      \textbf{few} & 可数名词复数 & \textbf{几乎没有}（强烈否定意味，如 \textit{few authors}） \\
      \textbf{a few} & 可数名词复数 & \textbf{有几个，少数几个}（肯定意味，如 \textit{a few celebrities}） \\
      \textbf{little} & 不可数名词 & \textbf{几乎没有}（强烈否定意味） \\
      \textbf{a little} & 不可数名词 & \textbf{有一点儿}（肯定意味） \\
      \bottomrule
    \end{tabularx}
  \end{table}
  \textbf{3. 省略句逻辑还原}：原文中的 \textit{not interpersonal, influence} 是高度紧凑的省略写法，其完整逻辑是 \textit{not a function of interpersonal influence}，即“主要是媒体效应，而非人际影响力”。
\end{proposition}

\subsubsection*{$\blacktriangleright$ 重点词汇与音标}
\noindent
\begin{tabularx}{\textwidth}{l X p{7cm}}
  \toprule
  \textbf{单词} & \textbf{音标与词性} & \textbf{核心释义} \\
  \midrule
  \retrokey{observe} & {\citshape /əbˈzɜːv/} \textit{v.} & 观察，遵守，注意到，发表评论 \\
  \midrule
  \retrokey{due to} & {\citshape /djuː tuː/} \textit{prep.} & 由于，因为（等价于 \textbf{because of}） \\
  \midrule
  \retrokey{simply} & {\citshape /ˈsɪmpli/} \textit{adv.} & 仅仅，简单地；简直，确实（用于加强语气） \\
  \midrule
  \retrokey{presence} & {\citshape /ˈprezns/} \textit{n.} & 存在感，在场，影响力（\textbf{outsized presence} = 超大的影响力） \\
  \bottomrule
\end{tabularx}

\subsubsection*{$\blacktriangleright$ 翻译与真题启示}
\begin{remark}
  \textbf{参考译文}：研究者们的观点源于对社会影响力的一次简单调查：除了少数几个像奥普拉·温弗莱这样的名人（其超大的影响力主要是媒体效应，而非人际影响力）之外，即使是人群中最有影响力的个体也根本不会与那么多人发生互动。
\end{remark}

\par\vspace{10pt}\hrule\vspace{10pt}

% =============================================================================
\section{第 11 句：会计准则对第三方资产评估的规定 (2010-4)}

\begin{note}
  These rules say they must value some assets at the price a third party would pay, not the price managers and regulators would like them to fetch. 【27 词】
\end{note}

\subsubsection*{$\blacktriangleright$ 句法层级拆解}
\noindent
\begin{tabularx}{\textwidth}{p{3.2cm} X X}
  \toprule
  \textbf{语法成分} & \textbf{具体英文内容} & \textbf{解析与修饰关系} \\
  \midrule
  \textbf{主句主干} & These rules say & 这些规则规定说（主谓结构） \\
  \midrule
  \textbf{宾语从句} & (that) they must value some assets at the price… not the price… & 宾语从句。省略了引导词 that，从句中 value 为及物动词 \\
  \midrule
  \textbf{从句并列结构} & at the price [a third party would pay], not [at] the price [managers... would like...] & 按照第三方愿意支付的价格评估，而非按照管理者想要的价格评估（并列省略了第二个 at） \\
  \midrule
  \textbf{定语从句 1} & a third party would pay & 修饰第一个 price（省略了关系代词 that/which，在从句中作 pay 的宾语） \\
  \midrule
  \textbf{定语从句 2} & managers and regulators would like them to fetch & 修饰第二个 price（同样省略了关系代词，them 指代 some assets） \\
  \bottomrule
\end{tabularx}

\begin{proposition}[连词的并列省略与定语从句省略]
  \textbf{1. 常见的省略规则 $\rightarrow$ 前提是绝不产生语义歧义}：
  \begin{itemize}
    \item \textbf{宾语从句引导词} $\rightarrow$ 当 \textbf{that} 引导宾语从句时，常被直接省略（如本句中的 \textit{say (that) they must...}）。
    \item \textbf{定语从句关系代词} $\rightarrow$ 当关系代词在定语从句中充当**宾语**时，可直接省略（如 \textit{the price (that) a third party would pay}）。
  \end{itemize}
  \textbf{2. 并列结构中的介词省略}：原文中的 \textit{at the price ... not the price ...}，在后半部分并列时直接省略了重复的介词 \textbf{at}。还原后为 \textit{value some assets at the price ... not at the price...}。
\end{proposition}

\subsubsection*{$\blacktriangleright$ 重点词汇与音标}
\noindent
\begin{tabularx}{\textwidth}{l X p{7cm}}
  \toprule
  \textbf{单词} & \textbf{音标与词性} & \textbf{核心释义} \\
  \midrule
  \retrokey{deem} & {\citshape /diːm/} \textit{v.} & 认为，视为（等同于 \textbf{consider}） \\
  \midrule
  \retrokey{appraise} & {\citshape /əˈpreɪz/} \textit{v.} & 估量，评估，鉴定（如评估资产价值） \\
  \midrule
  \retrokey{fetch} & {\citshape /fetʃ/} \textit{v.} & 售得，卖得（某价钱）；去取来（如 \textit{fetch some assets}） \\
  \midrule
  \retrokey{bring in} & {\citshape /brɪŋ ɪn/} \textit{phr.} & 带来（收益），年收入达……（如 \textit{bring in \$20,000}） \\
  \bottomrule
\end{tabularx}

\subsubsection*{$\blacktriangleright$ 翻译与真题启示}
\begin{remark}
  \textbf{参考译文}：这些规则规定，他们必须按照第三方愿意支付的价格来评估某些资产，而非管理者和监管者希望它们能卖到的价格。
\end{remark}

\par\vspace{10pt}\hrule\vspace{10pt}

% =============================================================================
\section{第 12 句：国际会计标准委员会规则重组压力 (2010-4)}

\begin{note}
  The IASB says it does not want to act without overall planning, but the pressure to fold when it completes its reconstruction of rules later this year is strong. 【29 词】
\end{note}

\subsubsection*{$\blacktriangleright$ 句法层级拆解}
\noindent
\begin{tabularx}{\textwidth}{p{3.2cm} X X}
  \toprule
  \textbf{语法成分} & \textbf{具体英文内容} & \textbf{解析与修饰关系} \\
  \midrule
  \textbf{第一分句主干} & The IASB says it does not want to act without overall planning & IASB表示其不想在缺乏总体规划时采取行动（转折前分句） \\
  \midrule
  \textbf{转折连词} & but & 但是，然而（强并列转折） \\
  \midrule
  \textbf{第二分句主语} & the pressure to fold & 屈服的压力（to fold 不定式作后置定语，修饰 pressure） \\
  \midrule
  \textbf{时间状语从句} & when it completes its reconstruction of rules later this year & 当今年晚些时候其完成规则重组时（when 引导时间状语从句，插入在主谓之间） \\
  \midrule
  \textbf{第二分句谓语} & is strong & 是巨大的，是强烈的（主谓一致） \\
  \bottomrule
\end{tabularx}

\begin{proposition}[高阶学术英语的翻译微调与重组]
  直译某些学术名词会显得生硬和空洞。在汉译时，应进行语境化“标志性”微调：
  \begin{itemize}
    \item \textbf{pressure to fold} $\rightarrow$ 直译为“折叠的压力” $\rightarrow$ 优化为“\retrokey{屈服的压力 / 妥协的压力}”。
    \item \textbf{reconstruction of rules} $\rightarrow$ 直译为“规则的重建” $\rightarrow$ 优化为“\retrokey{会计准则的修订与重组}”。
    \item \textbf{the pressure is strong} $\rightarrow$ 直译为“压力是强大的” $\rightarrow$ 优化为“\retrokey{面临巨大的妥协压力}”。
  \end{itemize}
\end{proposition}

\subsubsection*{$\blacktriangleright$ 重点词汇与音标}
\noindent
\begin{tabularx}{\textwidth}{l X p{7cm}}
  \toprule
  \textbf{单词} & \textbf{音标与词性} & \textbf{核心释义} \\
  \midrule
  \retrokey{overall} & {\citshape /ˌəʊvərˈɔːl/} \textit{adj./adv.} & 全面的，总体的；大体上 \\
  \midrule
  \retrokey{uniform} & {\citshape /ˈjuːnɪfɔːm/} \textit{adj./n.} & 统一的，一致的；制服 \\
  \midrule
  \retrokey{fold} & {\citshape /fəʊld/} \textit{v./n.} & 屈服，折叠，倒闭（\textbf{pressure to fold} = 屈服压力） \\
  \midrule
  \retrokey{reconstruction} & {\citshape /ˌriːkənˈstrʌkʃn/} \textit{n.} & 重建，改造，重组 \\
  \bottomrule
\end{tabularx}

\subsubsection*{$\blacktriangleright$ 翻译与真题启示}
\begin{remark}
  \textbf{参考译文}：IASB（国际会计标准委员会）表示，不想在没有总体规划的情况下采取行动。但在今年晚些时候完成其制度修订时，他们将面临巨大的妥协压力。
\end{remark}

\par\vspace{10pt}\hrule\vspace{10pt}

% =============================================================================
\section{第 13 句：麦克里维警告IASB政治现实真空 (2010-4)}

\begin{note}
  Charlie McCreevy, a European commissioner, warned the IASB that it did "not live in a political vacuum" but "in the real world" and that Europe could yet develop different rules. 【30 词】
\end{note}

\subsubsection*{$\blacktriangleright$ 句法层级拆解}
\noindent
\begin{tabularx}{\textwidth}{p{3.2cm} X X}
  \toprule
  \textbf{语法成分} & \textbf{具体英文内容} & \textbf{解析与修饰关系} \\
  \midrule
  \textbf{主句主语} & Charlie McCreevy, a European commissioner & 欧盟委员查理·麦克里维（commissioner 为同位语插入） \\
  \midrule
  \textbf{谓语与宾语} & warned the IASB & 警告IASB（双宾结构的前置间接宾语） \\
  \midrule
  \textbf{宾语从句 1} & that it did "not live in a political vacuum" but "in the real world" & 警告说其并非生活在政治真空而是现实世界（not... but... 介词并列结构） \\
  \midrule
  \textbf{宾语从句 2} & (and) that Europe could yet develop different rules & 并且警告说欧洲仍然有可能制定出不同的规则（and 连接的并列宾语从句） \\
  \bottomrule
\end{tabularx}

\begin{proposition}[“warn”经典句型与“vac-”词根辨析]
  \textbf{1. warn 的高频句型}：
  \begin{itemize}
    \item \textbf{warn sb. to do sth.} $\rightarrow$ 警告某人去做某事（不定式作宾补）。
    \item \textbf{warn sb. that...} $\rightarrow$ 警告某人说……（warn 后直接连接双宾语从句）。
  \end{itemize}
  \textbf{2. 词根 vac-（空）深度联想}：
  \begin{itemize}
    \item \textbf{vacuum} (\textit{n.} 真空 $\rightarrow$ 引申为“政治真空”) \quad \textbf{vanish} (\textit{v.} 突然消失 $\rightarrow$ \textit{vani-}变空，没了)
    \item \textbf{vacant} (\textit{adj.} 空缺的，空洞的) \quad \textbf{vanity} (\textit{n.} 虚荣，空虚)
  \end{itemize}
  \textbf{3. could yet 的强调含义}：\textbf{could yet} 表明“仍然存在发生的可能性”，译为“仍然可能、还是有可能会……”，表示一种保留的强硬态度。
\end{proposition}

\subsubsection*{$\blacktriangleright$ 重点词汇与音标}
\noindent
\begin{tabularx}{\textwidth}{l X p{7cm}}
  \toprule
  \textbf{单词} & \textbf{音标与词性} & \textbf{核心释义} \\
  \midrule
  \retrokey{commissioner} & {\citshape /kəˈmɪʃənə/} \textit{n.} & 委员，专员（如 \textit{European commissioner} 欧盟专员） \\
  \midrule
  \retrokey{foretell} & {\citshape /fɔːˈtel/} \textit{v.} & 预告，预言（等同于 \textbf{predict}） \\
  \midrule
  \retrokey{vacuum} & {\citshape /ˈvækjuəm/} \textit{n.} & 真空，真空吸尘器；政治真空，空白 \\
  \midrule
  \retrokey{vanish} & {\citshape /ˈvænɪʃ/} \textit{v.} & 突然消失，销声匿迹 \\
  \bottomrule
\end{tabularx}

\subsubsection*{$\blacktriangleright$ 翻译与真题启示}
\begin{remark}
  \textbf{参考译文}：欧洲委员查理·麦克里维警告IASB，它并非生活在“政治真空”中，而是生活在“现实世界”里，欧洲仍然可能制定出不同的规则。
\end{remark}

\par\vspace{10pt}\hrule\vspace{10pt}

% =============================================================================
\section{第 14 句：萧条市场与银行因恐亏损瘫痪 (2010-4)}

\begin{note}
  And dead markets partly reflect the paralysis of banks which will not sell assets for fear of booking losses, yet are reluctant to buy all those supposed bargains. 【27 词】
\end{note}

\subsubsection*{$\blacktriangleright$ 句法层级拆解}
\noindent
\begin{tabularx}{\textwidth}{p{3.2cm} X X}
  \toprule
  \textbf{语法成分} & \textbf{具体英文内容} & \textbf{解析与修饰关系} \\
  \midrule
  \textbf{主句主干} & And dead markets partly reflect the paralysis of banks & 萧条的市场在一定程度上反映了银行的瘫痪状况（主谓宾结构） \\
  \midrule
  \textbf{定语从句 1} & which will not sell assets for fear of booking losses & 修饰 banks（which 引导定语从句，在从句中作主语，修饰银行不愿售出资产的行为） \\
  \midrule
  \textbf{原因状语} & for fear of booking losses & 由于担心账面出现损失（作定语从句内部的原因状语，表示心理担忧） \\
  \midrule
  \textbf{定语从句 2} & yet (which) are reluctant to buy all those supposed bargains & 但它们也不愿购买那些所谓的廉价资产（由 yet 连接的第二个并列定语从句，省略了主语 which） \\
  \bottomrule
\end{tabularx}

\begin{proposition}[连词 yet 的多维用法与介词短语定状判定]
  \textbf{1. 连词 yet 与副词 yet 的区别判定}：
  \begin{table}[H]
    \centering
    \small
    \begin{tabularx}{\textwidth}{p{2.5cm} p{2.8cm} X}
      \toprule
      \textbf{词性} & \textbf{所处语境} & \textbf{核心释义与用法} \\
      \midrule
      \textbf{副词} & 否定句句尾 & 译为“\textbf{还（没有）}”（如 \textit{I haven't received it yet.}） \\
      \textbf{副词} & 肯定句中 & 译为“\textbf{仍然，还}”（如 \textit{He is yet a child.}） \\
      \textbf{连词} & 连接并列分句 & 译为“\textbf{但是，然而}”（本句中 \textit{yet are reluctant to...}） \\
      \bottomrule
    \end{tabularx}
  \end{table}
  \textbf{2. 介词短语作“定语”还是“状语”的判定法则}：\\
  判定规则 $\rightarrow$ 检查介词短语前是否有可修饰的**特定名词**。如有且语义通顺则为定语；若无法匹配，则多为修饰全句的状语。\\
  本句中的 \textit{for fear of...} 前虽有名词 \textit{assets}，但如果将其作为定语修饰 \textit{assets}（译作“以免担心损失的资产”），语义极其荒谬 $\rightarrow$ 故判定为**原因状语**（译作“因为担心损失而不愿出售资产”）。
\end{proposition}

\subsubsection*{$\blacktriangleright$ 重点词汇与音标}
\noindent
\begin{tabularx}{\textwidth}{l X p{7cm}}
  \toprule
  \textbf{单词} & \textbf{音标与词性} & \textbf{核心释义} \\
  \midrule
  \retrokey{paralysis} & {\citshape /pəˈræləsɪs/} \textit{n.} & 瘫痪，麻痹，无能为力（区别于 \textbf{paralyze} v.） \\
  \midrule
  \retrokey{reluctant} & {\citshape /rɪˈlʌktənt/} \textit{adj.} & 不情愿的，勉强的（\textbf{be reluctant to do} 不愿做……） \\
  \midrule
  \retrokey{for fear of} & {\citshape /fɔː fɪər ɒv/} \textit{phr.} & 唯恐，以免，因为担心……（常作原因状语） \\
  \midrule
  \retrokey{bargain} & {\citshape /ˈbɑːɡən/} \textit{n./v.} & 特价商品，廉价资产；讨价还价 \\
  \bottomrule
\end{tabularx}

\subsubsection*{$\blacktriangleright$ 翻译与真题启示}
\begin{remark}
  \textbf{参考译文}：萧条的市场在一定程度上反映了银行的瘫痪状况，因为他们为了避免账面损失不愿出售资产，但也不愿购买那些所谓的廉价资产。
\end{remark}

\par\vspace{10pt}\hrule\vspace{10pt}

% =============================================================================
\section{第 15 句：FASB与IASB不顾集团反对整顿期权 (2010-4)}

\begin{note}
  The FASB and IASB have been exactly that, cleaning up rules on stock options and pensions, for example, against hostility from special interests. 【21 词】
\end{note}

\subsubsection*{$\blacktriangleright$ 句法层级拆解}
\noindent
\begin{tabularx}{\textwidth}{p{3.2cm} X X}
  \toprule
  \textbf{语法成分} & \textbf{具体英文内容} & \textbf{解析与修饰关系} \\
  \midrule
  \textbf{主句主干} & The FASB and IASB have been exactly that & FASB与IASB的确一向如此（that 指代上文提到的 independent and combative 独立且强势的） \\
  \midrule
  \textbf{伴随状语 (分词)} & cleaning up rules on stock options and pensions & 整理并整顿关于股票期权和养老金的制度（现在分词短语作伴随状语） \\
  \midrule
  \textbf{插入语} & for example & 例如（双逗号隔开，移除后不影响句子完整） \\
  \midrule
  \textbf{伴随状语 (介词)} & against hostility from special interests & 不顾来自特殊利益集团的反对与敌意（介词短语作伴随状语） \\
  \bottomrule
\end{tabularx}

\begin{proposition}[伴随状语的五大语法形式与翻译翻译]
  \textbf{1. 伴随状语的五大核心表现形式}（表示伴随谓语同时发生的状态）：
  \begin{itemize}
    \item \textbf{分词短语} $\rightarrow$ \textit{He stood there, surrounded by reporters.} (他站在那里，被记者包围)。
    \item \textbf{形容词短语} $\rightarrow$ \textit{He sat on the sofa, silent and motionless.} (他坐在沙发上，一言不发)。
    \item \textbf{介词短语} $\rightarrow$ \textit{They sat by the fireplace, in complete silence.} (他们沉默地坐在壁炉旁)。
    \item \textbf{独立主格} $\rightarrow$ \textit{He ran to the station, his bag swinging behind him.} (他跑向车站，包包晃荡)。
    \item \textbf{with 复合结构} $\rightarrow$ \textit{He ran into the room, with his dog barking loudly.} (他跑进房间，狗大叫)。
  \end{itemize}
  \textbf{2. 翻译语境化处理}：原文中的短语 \textbf{cleaning up} 字面义为“清理、打扫”。但因其动作承受对象是 \textit{rules}（制度规章），直接翻译会显得生硬，在汉译时应将其“语意优化”翻译为“**整顿**制度 / **规范**制度”，更加精准地道。
\end{proposition}

\subsubsection*{$\blacktriangleright$ 重点词汇与音标}
\noindent
\begin{tabularx}{\textwidth}{l X p{7cm}}
  \toprule
  \textbf{单词} & \textbf{音标与词性} & \textbf{核心释义} \\
  \midrule
  \retrokey{combative} & {\citshape /kəmˈbætɪv/} \textit{adj.} & 好斗的，强势的，好争论的 \\
  \midrule
  \retrokey{pension} & {\citshape /ˈpenʃn/} \textit{n./v.} & 养老金，退休金；发给……养老金 \\
  \midrule
  \retrokey{hostility} & {\citshape /hɒˈstɪləti/} \textit{n.} & 敌意，敌对，反对，抵制 \\
  \midrule
  \retrokey{special interests} & {\citshape /ˌspeʃl ˈɪntrəsts/} \textit{n.} & 特殊利益集团，特权集团 \\
  \bottomrule
\end{tabularx}

\subsubsection*{$\blacktriangleright$ 翻译与真题启示}
\begin{remark}
  \textbf{参考译文}：美国财务会计标准委员会和国际会计标准委员会的确一向如此独立和强势。例如，他们不顾某些特殊利益集团的反对，整顿股票期权和养老金制度。
\end{remark}

\par\vspace{10pt}\hrule\vspace{10pt}

% =============================================================================
\section{第 16 句：丰富多样刺激因素环境对智力的开发 (1990-2)}

\begin{note}
  The child who is raised in an environment where there are many stimuli which develop capacity for appropriate responses will experience greater intellectual development. 【29 词】
\end{note}

\subsubsection*{$\blacktriangleright$ 句法层级拆解}
\noindent
\begin{tabularx}{\textwidth}{p{3.2cm} X X}
  \toprule
  \textbf{语法成分} & \textbf{具体英文内容} & \textbf{解析与修饰关系} \\
  \midrule
  \textbf{主句主干} & The child [...] will experience greater intellectual development. & 孩子可能会经历更好的智力发育（主谓宾结构） \\
  \midrule
  \textbf{嵌套定从 1} & who is raised in an environment [...] & 修饰 The child（第一层定从，who 作主语） \\
  \midrule
  \textbf{嵌套定从 2} & where there are many stimuli [...] & 修饰 environment（嵌套于第1层内的第2层定从，where 作地点状语） \\
  \midrule
  \textbf{嵌套定从 3} & which develop capacity for appropriate responses & 修饰 stimuli（嵌套于第2层内的第3层定从，which 作主语，指代刺激因素的能力） \\
  \bottomrule
\end{tabularx}

\begin{proposition}[多层从句嵌套判定与情态动词的共性用法]
  \textbf{1. 长难句的多层从句嵌套拆解法则}：\\
  长难句最核心的设障手段是多层定从嵌套。阅读时，必须先理清各层关系词的指代对象：
  \begin{itemize}
    \item \textbf{who} $\rightarrow$ 指代 \textit{The child} $\rightarrow$ 孩子是在怎样的环境中长大的？
    \item \textbf{where} $\rightarrow$ 指代 \textit{an environment} $\rightarrow$ 环境中包含着什么？（包含着丰富的刺激因素）
    \item \textbf{which} $\rightarrow$ 指代 \textit{stimuli} $\rightarrow$ 这些刺激因素起到了什么作用？（开发了做出合适反应的能力）
  \end{itemize}
  \textbf{2. 动词词根 cap-（抓、拿、头）高频衍生链}：
  \begin{itemize}
    \item \textbf{capacity} (\textit{n.} 能力，容量) \quad \textbf{capture} (\textit{v.} 俘获，捕获)
    \item \textbf{capacious} (\textit{adj.} 宽敞的，容量大的) \quad \textbf{capital} (\textit{n./adj.} 首都，资本；大写的)
  \end{itemize}
  \textbf{3. 情态动词 will 的“共性”特殊翻译}：原文中的 \textbf{will experience} 并非特指将来某个特定时间点发生的动作，而是指科学上的**共性规律或客观趋势**。在翻译时，应当将其委婉地译为“**可能会**经历 / **往往会**得到……”，更符合汉语的逻辑严密性。
\end{proposition}

\subsubsection*{$\blacktriangleright$ 重点词汇与音标}
\noindent
\begin{tabularx}{\textwidth}{l X p{7cm}}
  \toprule
  \textbf{单词} & \textbf{音标与词性} & \textbf{核心释义} \\
  \midrule
  \retrokey{raise} & {\citshape /reɪz/} \textit{v.} & 抚养；筹集（资金）；招聘；提出（及物动词，区别于 \textbf{rise} 不及物） \\
  \midrule
  \retrokey{stimuli} & {\citshape /ˈstɪmjulaɪ/} \textit{n.} & 刺激，促进因素（复数形式，单数为 \textbf{stimulus}） \\
  \midrule
  \retrokey{capacity} & {\citshape /kəˈpæsəti/} \textit{n.} & 能力，接受力，容积，容量 \\
  \midrule
  \retrokey{appropriate} & {\citshape /əˈprəʊpriət/} \textit{adj./v.} & 合适的，恰当的；拨出（款项），挪用 \\
  \midrule
  \retrokey{intellectual} & {\citshape /ˌɪntəˈlektʃuəl/} \textit{adj./n.} & 智力的，理智的；脑力劳动者，知识分子 \\
  \bottomrule
\end{tabularx}

\subsubsection*{$\blacktriangleright$ 翻译与真题启示}
\begin{remark}
  \textbf{参考译文}：一个在具有丰富多样刺激因素的环境中长大的孩子，其智力水平可能会更高，因为这些刺激因素能培养孩子灵活应变的能力。
\end{remark}

\par\vspace{20pt}\hrule\vspace{20pt}

% =============================================================================
% 附录：常用语法框架速查
% =============================================================================
\appendix
\chapter{常用语法框架速查}

为了便于您在自测和做真题时快速核对，以下为您总结了长难句中最核心的三大句法结构。

\section{定语从句引导词速查}
\noindent
\begin{tabularx}{\textwidth}{l X X}
  \toprule
  \textbf{先行词类型} & \textbf{从句主语位置（不可省略）} & \textbf{从句宾语位置（可以省略）} \\
  \midrule
  \textbf{人} & who & whom / who \\
  \midrule
  \textbf{物} & which & which \\
  \midrule
  \textbf{人或物} & that & that \\
  \midrule
  \textbf{时间名词} & when (作状语) & that/which (作宾语时可省略) \\
  \midrule
  \textbf{地点名词} & where (作状语) & that/which (作宾语时可省略) \\
  \midrule
  \textbf{原因名词 (reason)} & why (作状语) & that/which (作宾语时可省略) \\
  \bottomrule
\end{tabularx}

\section{被动句汉译三大核心策略}
\begin{enumerate}
  \item \textbf{增加泛指主语} $\rightarrow$ 当原文施事者不明时，在汉语中增加“人们、我们、大家”作主动主语。
  \item \textbf{施事者明确时直接移前} $\rightarrow$ 原文出现 \textit{by...} 结构时，直接把 \textit{by} 后的施事者移到动词前作主语。
  \item \textbf{化被动为主动，调整语序} $\rightarrow$ 很多英语中的被动句在汉语中可以直接用主动句式自然流露（如“可以找到”代替“被找到”）。
\end{enumerate}

\section{分词短语作状语的类型与判定}
\noindent
\begin{tabularx}{\textwidth}{l p{5cm} X}
  \toprule
  \textbf{分词类型} & \textbf{经典句型示例} & \textbf{语法功能与逻辑关系} \\
  \midrule
  \textbf{现在分词 (表主动)} & approving a patent... & 伴随状语（译为：同时批准了……） \\
  \midrule
  \textbf{过去分词 (表被动)} & published in 1998... & 后置定语或时间/原因状语（译为：已出版的……） \\
  \midrule
  \textbf{独立主格结构} & the work done, he left... & 独立的时间或伴随状语 \\
  \bottomrule
\end{tabularx}

\end{document}
